% \begin{abstract}
% \YJ{TODO}
% Dynamic graphs demand both high-throughput structural updates and efficient, repeated
% downstream analytics. We present {\name}, a query-agnostic Graph Intermediate Representation
% that maintains multi-granular connectivity invariants—weakly connected components
% (WCC), strongly connected components (SCC) with a condensation DAG, and biconnected
% components (BiCC) organized as block-cut (BC) trees—and updates them incrementally
% using Scope–Prune–Update and Filter–Merge–Update strategies. By
% pruning at the DAG level and exposing independent BiCC blocks for concurrent
% execution, GIR enables block-level task parallelism that substantially reduces
% redundant downstream recomputation. We implement a GIR-enabled execution framework
% and evaluate it on seven real-world graphs (up to large scale), achieving up to
% 3.0× speedup on BFS, 2.7× on SSSP, and 1.7× on betweenness centrality while
% maintaining competitive update throughput and strong multi-thread scalability.
% \YJ{TODO}
% \end{abstract}
Dynamic graphs demand both high-throughput structural updates and efficient downstream analytics. 
Existing methods are either \textit{history-blind} as they discard prior computation after each update, incurring large repeated overhead
or \textit{query-dependent} as they reuse results but are tied to specific algorithms or parameters, offering no general
reuse. 
Neither approach captures algorithm-independent structural invariants as a persistent, reusable layer.
Inspired by compiler intermediate representations, we propose \textbf{GIR}, a
query-agnostic \emph{Graph Intermediate Representation} that decouples dynamic
graph updates from downstream query execution. GIR maintains a three-level
structural hierarchy of weakly connected components, strongly connected
components with a condensation DAG, and biconnected components
organized as block-cut trees. To keep GIR current under continuous updates,
we design strategies to confine each repair to the minimum affected region. GIR further enables a two-level parallel execution model that combines task parallelism across
independent blocks with data parallelism within each block. 
% Evaluated on
% seven real-world graphs, GIR achieves up to $3.0\times$ speedup on BFS,
% $2.7\times$ on SSSP, and $1.7\times$ on betweenness centrality over
% state-of-the-art systems, while maintaining competitive update throughput and
% near-linear scalability.
% We design {GIR}, a query-agnostic graph intermediate representation
% that bridges high-throughput dynamic graph updates and efficient downstream
% analytics. 
% GIR persistently maintains a three-level structural hierarchy of WCC, SCC, and BiCC, confines maintenance to the minimum affected region, and enables task parallelism.
% SCC with a condensation DAG, and BiCC with BC trees. 
% We confine maintenance to the minimum affected region, avoiding global recomputation.
% GIR further enables task parallelism.
Experiments on seven real-world graphs show that GIR achieves 1.8$\times$ more  update throughput than state-of-the-art methods.
For graph analytics, GIR achieves an average of 1.2$\times$, 1.2$\times$, and 1.7$\times$ speedups for breadth-first search, single-source shortest path, and betweenness centrality, respectively.